% !TeX root = ../papers/nc-p3.tex
\section{三次スクロール型の中心生成元・最小関係と短い正則性検査}
\label{sec:scroll-spin-center-generators}
\begin{theorem}\label{thm:scroll-spin-center-generators}
前節第四型の中心の最小生成元は圧縮次数 $1,2,3$ にそれぞれ $5,3,4$ 個あり、
最小関係は次数 $2,3,4,5,6$ にそれぞれ $3,6,10,6,10$ 個ある。
一般の次数 $1,2$ の二元を割ると、残る環は二つの分岐点を印した種数9の超楕円曲線の
半標準環である。さらに次数2の二元を割ったArtin環は分岐点の残る十八点の位置によらず同型である。
\end{theorem}
\begin{proof}
一般八次曲面から一般標準曲線 $C$ を切る。両元は正則列で、$C$ に沿って中心の次数安定化がなく、
切断列と中間コホモロジーの消滅により中心の商は $\bigoplus_n H^0(C,n\theta)$,
$\theta=3J+P+Q$ である。幾何学的一般体で曲線を表示する。
二重被覆の底の座標を $s,t$ とし、二分岐点を $st=0$ へ運ぶ。残る平方自由な分岐式を $r_{18}(s,t)$ とする。
補助環
\[
    k[s,t,b,c]/(b^2-st,\ c^2-r_{18}),\qquad
             \deg(s,t,b,c)=(1,1,1,9)
\]
で、重み $4n$ かつ $(b,c)\mapsto(-b,-c)$ の指標 $(-1)^n$ の部分が半標準環の次数 $n$ である。
実際 $S_j=k[s,t]_j$ とおくと、$n=2a$ では $S_{8a}\oplus bcS_{8a-10}$、
$n=2a+1$ では $bS_{8a+3}\oplus cS_{8a-5}$ になる。偶数側は二重被覆式
$y^2=st r_{18}$ の $y=bc$、奇数側は二分岐点で一位の極を許した反不変部 $y/(st)$ を
固定切断 $b$ と掛けたものである。従って積もこの表示に一致する。

最小生成元は $x_i=b s^{3-i}t^i$, $y_i=c s^{3-i}t^i$ ($0\le i\le3$) と $u=s^8,v=t^8$ で、
半標準次数は順に $1,3,2$ である。二次の $x_ix_j$ は $stS_6$ を張り、$u,v$ を加え全八次になる。
三次では $cS_3$ が初めて現れる。それ以降は八次式の積と、偶数の反不変部に対する $x_i y_j$ で全てを張る。
二元 $u,v$ は正則列で、その商は補助環で $s^8=t^8=0$ と置いて同じ重み・指標を抜いた部分である。
実際補助環は $k[s^8,t^8]$ 上で、指数 $0\le i,j<8$ の $s^it^j$ と $1,b,c,bc$ の積を
基底とする自由加群であり、重み剰余と指標で分かれる。商では $r_{18}$ は自動的に零になる。

この商の関係を数える。$x,y$ の個数が固定された単項式で添字の総和を $j$ とする。
$x$ のみの二因子では同じ $j$ を同一視する三本、$x,y$ 一つずつでは同じ $j$ を同一視する九本が要る。
これらの関係のもと同じ種類数と $j$ を持つ単項式は、$y$ が高々一個なら全て等しい。
二因子間で添字を一単位ずつ移して揃えられるからである。
$x$ 三つでは $3\le j\le6$ の四類だけが生きるので新たに六本、
四つでは三次関係からまだ消えない $j=6$ の一類を消し、それ以上は全て零である。
$x$ 二つと $y$ 一つでも $3\le j\le6$ だけが生き、新たに六本が必要である。
$x$ 三つと $y$ 一つでは $j=6$ のみが残り、他は既にこれらの関係で消える。
最後に $y_i y_j=0$ が対称十本あり、二つ以上の $y$ を含む全てが零になる。
これで残る基底は上の打切り単項式に一対一で、重複を除いて示した本数は最小である。

正則列を割っても最小関係の次数が戻ることを述べる。パラメータは経路ではなく可換中心の
最小多項式表示の変数に取れる。最初の二元はそれぞれ一般一次式と、三つの新二次元の一つを含む
一般二次式であり、最後の $u,v$ は明示した残りの新二次元である。
各商の生成元を持ち上げれば次数帰納で全環を生成する。多項式表示上の最小自由分解を
各変数で順に割ると、商環での正則性により高次Torがなく最小性も保たれる。
従って求めたBetti数が元の中心のものとなる。体拡大はこの次数別の数を変えない。
\end{proof}

\begin{proposition}\label{prop:central-arrow-module-cutoff}
定理\ref{thm:successively-central-cutoff}で四元が初めから中心にあるとき、経路次元の検査上限
$3b$ を $\max(L+2,d_1,\ldots,d_4)$ に下げられる。剰余の上限 $L+1$、複体の上限 $L+4$ はそのままでよい。
十五次元閉包の有限四辺で正則整域となることは、経路長 $(4,4,8,8)$ の中心元でこの検査を通ることと同値である。
従って経路次元は22次、中心剰余は21次、標準複体は24次までで足りる。
基点が高々有限な型では $(4,4,4,8)$ が使え、それぞれ18、17、20次で足りる。
\end{proposition}
\begin{proof}
四変数の重み付き多項式環 $P$ 上、最終剰余の斉次経路基底 $v$ の各持ち上げに一つずつ
自由記号を置いた加群 $F$ を作る。頂点ごとに単位も取る。中心性と次数帰納から $F\to B$
は全射、標準経路次元を仮定した範囲で同型である。$F$ 上の各右矢作用をその自由基底での
一意な経路展開で定め、$P$ 線型に延長する。$v$ に二次関係を作用させたものは $B$ で消え、
その長さは $L+2$ 以下なので $F$ でも消える。従って真の右 $B$ 加群になり、先の射は $B$ 線型でもある。
単位記号に各基底経路を作用させるとその記号が戻り、各中心元の経路式を作用させると
各頂点で対応する多項式変数と単位の積が戻る。これらは検査済みの次数だからである。
作用の $P$ 線型性から全ての多項式倍も単位記号で生成される。従って頂点単位からの右正則加群の写像
$B\to F$ は全射で、合成 $B\to F\to B$ が恒等なので同型である。全次数で標準経路次元が得られ、
あとは元の有限検査の十分性を使える。

前節の四型には四次中心系の固定因子がない。完全交叉の三型では基点は空または有限、
第四型では一般八次切断上に二点のみなので同じく固定因子はない。
よって一般の二つの四次切断の共通台は射影的に高々曲線で、
基点のない八次切断を二つ追加すれば空になる。基点有限なら四次を三つ先に選べる。
逆に有限検査の成立は正則性と中心有限性を与え、既出のNoether窓の定理で合成整域になる。
\end{proof}
